% !TEX root = ../main.tex
\section{Simulation Testbed}
\label{sec:simulation-testbed}

\draftnote{TO BE REVIEWED}

% This section describes the physical testbed used to evaluate learned prosthetic
% reference trajectories and to validate the control strategy. 
% The testbed is necessary for the evaluation of the trajectory: the learned policy
% does not directly actuate muscles, or motor torques. Instead, it
% provides a prosthetic kinematic reference used by the control chain (cascade controller and motor drive) while the biological-side is driven by sperimental prescribed kinematic references. The simulator (based on OpenSim's \textit{Computed Muscles Control}) makes it possible by decoupling the prosthetic side from the biological one. This separation is
% central to the paper: a learned trajectory is meaningful only if it can be
% realized by a physical series-elastic prosthesis inside a muscle-driven gait
% simulation.
\noindent 
This section describes the physical testbed used both to evaluate learned
prosthetic reference trajectories and to validate the prosthetic control chain.
The learned policy does not directly command muscles or motor torques.
Instead, it provides a prosthetic kinematic reference that is tracked through the
controllers. On the biological side,
experimental kinematics are used as prescribed tracking references: the simulator (based on OpenSim's \textit{Computed Muscles Control}) computes the generalized force demand and distributes it through muscle recruitment and residual actuators. Therefore the prosthetic and biological coordinates are treated by two distinct control paths. The prosthetic side is driven by the
learned reference and SEA controller, whereas the biological side is tracked from
experimental kinematics through muscle-driven recruitment. Both sides remain
coupled in the same forward-dynamics OpenSim simulation. This separation is
central to the paper: a learned trajectory is meaningful only if it can be
realized by a physical series-elastic prosthesis inside a muscle-driven gait
simulation.

\subsection{Series-Elastic Prosthesis Model}
\label{subsec:sea-model}

\noindent
Series-elastic actuators consist of an electric motor connected to the joint through a torsional spring in series. Thus, the joint is not directly driven by the motor, but is actuated through the deflection of the spring. SEAs were selected to actuate the transfemoral prosthesis for two main reasons:
\begin{enumerate}
  \item The spring can store and release energy during the gait cycle. This may improve energy efficiency compared with conventional motor actuation and provides intrinsic compliance, allowing the actuator to attenuate impacts during gait.
  \item SEAs are mechanically simpler than variable stiffness actuators.
\end{enumerate}
\noindent
Each SEA was implemented as a custom OpenSim plugin. This choice preserves compatibility with the OpenSim software suite.

The SEA plant is described by two physical motor-side state variables:
\begin{equation}
  \theta_m,\quad \omega_m ,
\end{equation}
where \(\theta_m\) is the motor-side angle, \(\omega_m\) is the motor angular
velocity, joint coordinate angle is denoted by \(\theta_j\). The
elastic torque transmitted through the series spring is
\begin{equation}
  \tau_s = K(\theta_m - \theta_j),
  \label{eq:sea-spring}
\end{equation}
where \(K\) is the spring stiffness. 
\newline The SEA dynamics are
\begin{equation}
\dot{\theta}_m = \omega_m,\qquad
J_m\dot{\omega}_m = \tau_m - \tau_s - B_m\omega_m .
\label{eq:sea-dynamics}
\end{equation}
Here \(J_m\) is the motor inertia, \(B_m\) is the motor damping, and
\(\tau_m\) is the motor-side torque.
\newline 
Thus, the prosthetic joint is actuated by the torque transmitted through the
elastic element, rather than being directly driven by the motor. The resulting
torque-generation dynamics are therefore determined by the coupled behavior of
the motor and compliant transmission, with motor inertia, damping, and spring
stiffness governing the rate at which torque can be developed and delivered to
the joint.

\noindent The actuators are dimensioned as follows:

\begin{table}[!h]
  \centering
  \footnotesize
  \begin{tabular}{@{}lcccc@{}}
    \toprule
    SEA &
    \makecell{\(K\)\\\(\mathrm{N\,m/rad}\)} &
    \makecell{\(J_m\)\\\(\mathrm{kg\,m^2}\)} &
    \makecell{\(B_m\)\\\(\mathrm{N\,m\,s/rad}\)} &
    \makecell{\(F_{\mathrm{opt}}\)\\\(\mathrm{N\,m}\)} \\
    \midrule
    Knee  & 321 & 0.01 & 0.1 & 100 \\
    Ankle & 500 & 0.01 & 0.1 & 250 \\
    \bottomrule
  \end{tabular}
  \vspace{0.5em}
  \caption{SEA parameters adopted.}
  \label{tab:sea-plant-parameters}
\end{table}

Thus, the prosthetic joint is actuated by the torque transmitted through the elastic element, rather than being directly driven by the motor. The resulting torque-generation dynamics are therefore determined by the coupled behavior of the motor and compliant transmission, with motor inertia, damping, and spring stiffness governing the rate at which torque can be developed and delivered to the joint.

\paperplaceholder{Figure placeholder. Fig.~2 should show the two-body SEA
schematic with \(\theta_j\), \(\theta_m\), \(K\), \(B_m\), \(J_m\),
\(\tau_s\), and \(\tau_m\).}


\subsection{CMC-Like Muscle-Driven Simulation}
\label{subsec:cmc-like-simulation}

\noindent The prosthetic model has been evaluated using a custom simulator based on OpenSim's \textit{Computed Muscle Control} (CMC) tool. OpenSim Computed Muscle Control is a forward-dynamics tracking method that computes the muscle excitations required for a musculoskeletal model to follow a prescribed kinematic trajectory. At each time step, the simulated state is compared with the reference motion, desired accelerations are generated through a tracking controller, and muscle recruitment is solved to produce the generalized forces needed to realize those accelerations. Reserve actuators may be included as residual support when the modeled muscles cannot fully satisfy the required dynamics.
The limitation of the original CMC is that it does not allow to impose a custom control law on a subset of model coordinates.

\paperplaceholder{Figure placeholder. Fig.~3 should show the CMC loop}

The integration step is \(1\,\mathrm{ms}\), and the
validated forward-dynamics path uses the C++ SEA plugin with the
\texttt{rk4\_bypass} integration scheme.

The kinematic reference is read from the inverse-kinematics file and
preprocessed before it is used by the controller. The trajectory is resampled on
a \(1\,\mathrm{ms}\) grid and low-pass filtered with a zero-phase Butterworth
filter at \(6\,\mathrm{Hz}\). This preprocessing keeps the spline derivatives
usable for acceleration-level tracking and avoids injecting high-frequency IK
noise into the inverse-dynamics and static-optimization stages.

At each control instant, the CMC-like biological tracking loop computes desired
accelerations for the non-prosthetic coordinates:
\begin{equation}
\ddot{q}_{\mathrm{des}} =
\ddot{q}_{\mathrm{ref}}
+ K_{p,b}(q_{\mathrm{ref}} - q)
+ K_{d,b}(\dot{q}_{\mathrm{ref}} - \dot{q}).
\label{eq:bio-outer-loop}
\end{equation}
The default biological gains are \(K_{p,b}=100\) and \(K_{d,b}=20\), with
lower translation gains on the pelvis coordinates (\(25/10\)) to reduce
root-coordinate drift without forcing the muscle-driven joints through the
reserve actuators. The prosthetic coordinates are excluded from this biological
outer loop and from the biological static-optimization problem; they are handled
by the SEA controller.

The generalized force request is computed using a zero-actuator
inverse-dynamics projection. First, all muscle, reserve, and SEA controls are
removed. The baseline acceleration \(\ddot{q}_0\) is computed at the current
state with the actuator contributions zeroed. The acceleration deficit is then
projected through the configuration-dependent mass matrix:
\begin{equation}
\tau = M(q)\left(\ddot{q}_{\mathrm{des}} - \ddot{q}_0\right).
\label{eq:id-projection}
\end{equation}
The resulting generalized-force vector is split into a biological component and
a prosthetic component:
\begin{equation}
\tau_{\mathrm{bio}} = \tau[\mathcal{I}_{\mathrm{bio}}],
\qquad
\tau_{\mathrm{pros}} = \tau[\mathcal{I}_{\mathrm{pros}}].
\label{eq:force-split}
\end{equation}
Only \(\tau_{\mathrm{bio}}\) is passed to the muscle recruitment step. The
prosthetic component is retained as a diagnostic inverse-dynamics oracle, but it
is not used as a feed-forward command in the current prosthetic controller.

Biological recruitment is solved with a muscle-first static-optimization
problem. The optimizer distributes \(\tau_{\mathrm{bio}}\) across muscle
activations and reserve actuator controls:
\begin{equation}
\begin{aligned}
\min_{a,u_r}\quad
& \sum_i a_i^2 + w_r \sum_j u_{r,j}^2 \\
\mathrm{s.t.}\quad
& A_m(q,\dot{q},l_f)(a-a_{\min})
  + R_r(u_r \odot F_{\mathrm{opt},r})
  = \tau_{\mathrm{bio}}, \\
& a_{\min} \leq a_i \leq a_{\max}, \\
& -u_{r,\max} \leq u_{r,j} \leq u_{r,\max}.
\end{aligned}
\label{eq:static-optimization}
\end{equation}
The reserve weight is set to \(w_r=10^6\). Thus, reserves remain available to
close unmodeled or infeasible residual directions, but they are strongly
discouraged from becoming the primary tracking mechanism. In the analysis,
reserve usage is treated as a fidelity gauge rather than as a success metric:
low reserve usage indicates that the simulated motion is supported by the
model's muscles and prosthesis, whereas high reserve usage exposes a dynamic
inconsistency, a contact-model deficit, or an unmodeled actuation demand.

The CMC-like loop therefore provides two pieces of information that are
essential for evaluating learned prosthetic trajectories. First, it determines
whether the biological side can remain muscle-driven while the prosthesis
follows the generated reference. Second, it makes reserve and residual loads
available as quantitative diagnostics of physical validity. This prevents the
RL policy from being judged only by kinematic tracking or reward return.

\paperplaceholder{Figure placeholder. Fig.~1 should include the full testbed
flow: IK spline \(\rightarrow\) biological outer-loop accelerations
\(\rightarrow\) zero-actuator inverse dynamics \(\rightarrow\) static
optimization \(\rightarrow\) prosthetic cascade \(\rightarrow\) SEA plugin
\(\rightarrow\) OpenSim state update.}

\paperplaceholder{Table placeholder. Table~IV should summarize the AB06 model,
simulation window, time step, prosthetic coordinates, biological tracking gains,
static-optimization weights, and GRF mode.}

\subsection{Cascade/PI Prosthetic Reference Tracking}
\label{subsec:prosthetic-controller}

The two SEA coordinates are driven by a fixed high-level prosthetic tracking
controller implemented in Python. This controller receives the prosthetic
kinematic reference \((q_{\mathrm{ref}},\dot{q}_{\mathrm{ref}})\), compares it
with the current prosthetic coordinate state \((q,\dot{q})\), and computes the
normalized SEA command \(u\) sent to the C++ plugin. The current validated mode
is a position-to-velocity cascade followed by an inner velocity PI loop.

For each prosthetic coordinate, the outer proportional position loop computes a
desired cascade velocity:
\begin{equation}
\dot{q}_{\mathrm{cas}} =
\dot{q}_{\mathrm{ref}}
+ K_{p,o}(q_{\mathrm{ref}} - q).
\label{eq:cascade-velocity}
\end{equation}
The inner loop then converts the velocity tracking error into a torque command:
\begin{equation}
e_v = \dot{q}_{\mathrm{cas}} - \dot{q},
\label{eq:velocity-error}
\end{equation}
\begin{equation}
\tau_{\mathrm{cmd}} =
K_{p,i}e_v
+ K_{i,i}\int e_v\,dt.
\label{eq:cascade-torque}
\end{equation}
The integral contribution is clamped per joint to prevent wind-up. Finally, the
controller normalizes the requested torque by the corresponding SEA optimal
force:
\begin{equation}
u =
\operatorname{clip}\left(
\frac{\tau_{\mathrm{cmd}}}{F_{\mathrm{opt}}},
-1,1
\right).
\label{eq:sea-command}
\end{equation}
This \(u\) is the only prosthetic command injected into OpenSim. The C++ plugin
converts it into a desired spring torque,
\begin{equation}
\tau_{\mathrm{ref}} = F_{\mathrm{opt}}u.
\label{eq:spring-reference}
\end{equation}

Inside the plugin, the actuator drive closes a motor-side spring-torque servo.
Let \(\xi\) denote the integrated spring-torque tracking error. In the
non-impedance mode used by the validated model, the raw motor torque is
\begin{equation}
\tau_{m,\mathrm{raw}} =
\tau_{\mathrm{ref}}
+ K_p(\tau_{\mathrm{ref}}-\tau_s)
+ \operatorname{sat}(K_i\xi,\pm L_i)
- K_d\omega_m ,
\label{eq:motor-servo}
\end{equation}
with
\begin{equation}
\dot{\xi} = \tau_{\mathrm{ref}} - \tau_s.
\label{eq:integral-state}
\end{equation}
The integral contribution is limited to \(\pm L_i\), and the integrator is
frozen when either the integral contribution or the motor torque is saturated in
the direction of the error. The motor torque applied to the rotor is
\begin{equation}
\tau_m =
\operatorname{clip}(\tau_{m,\mathrm{raw}}, -500, 500)\,\mathrm{N\,m}.
\label{eq:motor-clamp}
\end{equation}
This \(\tau_m\) drives the SEA dynamics in Eq.~\eqref{eq:sea-dynamics}; the
joint still receives the spring torque \(\tau_s\) defined in
Eq.~\eqref{eq:sea-spring}.

The cascade gains used in the current configuration are reported in Table~III.
The knee uses \(K_{p,o}=18.85\,\mathrm{s^{-1}}\),
\(K_{p,i}=29.2\,\mathrm{N\,m\,s/rad}\),
\(K_{i,i}=1377\,\mathrm{N\,m/rad}\), and a \(50\,\mathrm{N\,m}\) inner
integral torque limit. The ankle uses
\(K_{p,o}=47.125\,\mathrm{s^{-1}}\),
\(K_{p,i}=2.8275\,\mathrm{N\,m\,s/rad}\),
\(K_{i,i}=213\,\mathrm{N\,m/rad}\), and a \(200\,\mathrm{N\,m}\) inner
integral torque limit. The embedded actuator-drive gains are
\((K_p,K_d,K_i)=(18,11,190)\) for the knee and \((11.3,11,123)\) for the
ankle, with \(L_i=100\,\mathrm{N\,m}\) for both joints.

This controller is deliberately not a learned low-level policy. It is a fixed,
validated tracking interface between the generated kinematic reference and the
SEA plant. Its role is to expose the correct failure mode to the trajectory
generator: when the reference is too aggressive, the controller cannot hide that
fact. It produces larger \(u\), larger motor torque demand, increased spring
tracking error, or saturation. Conversely, a feasible reference remains
trackable without changing the low-level SEA parameters. This is why the
reference governor introduced in Section~IV is part of the action interface
rather than a post-hoc smoothing operation.

\paperplaceholder{Table placeholder. Table~III should report cascade gains,
cascade integral limits, actuator-drive gains, actuator-drive integral and
motor-torque limits, reference-governor velocity/acceleration/jerk limits, and
the normalized command limits.}

\draftnote{Before moving this section to a final submission template, verify all
reported numerical gains against the final model/config snapshot used for the
experiments, and decide whether ``series-elastic actuator'' or ``SEA'' is used
as the dominant term.}
